\documentclass[11pt]{article}
\usepackage{amsmath,amssymb,geometry}
\geometry{margin=1in}
\newcommand{\Tr}{\operatorname{Tr}}
\newcommand{\HS}{\mathrm{HS}}
\begin{document}
\title{Eight symmetric qubits: second- and third-order anticoherent subspaces have maximum dimension one}
\author{Exact catalogue calculation}
\date{}
\maketitle

\paragraph{Claim.}
For pure symmetric eight-qubit subspaces,
\[
 k_{\max}(8,2)=k_{\max}(8,3)=1.
\]
The upper bound below applies to arbitrary complex subspaces. No choice of support, symmetry, real coefficients, or Majorana multiplicities is assumed.

\paragraph{1. A hypothetical two-dimensional subspace leaves a pure complement.}
The symmetric Hilbert space is the spin-$j=4$ representation, of dimension nine. Let $P$ be the orthogonal projector onto a hypothetical two-dimensional second-order anticoherent subspace. Polarization of the anticoherence conditions gives
\[
 PJ_aP=0,\qquad PJ_aJ_bP=\frac{j(j+1)}3\delta_{ab}P
 =\frac{20}{3}\delta_{ab}P,
 \qquad a,b\in\{x,y,z\}.
\]
For any orthonormal basis $u_1,u_2$ of its range, the eight vectors
\[
 u_r,\quad \sqrt{\frac3{20}}J_xu_r,\quad
 \sqrt{\frac3{20}}J_yu_r,\quad
 \sqrt{\frac3{20}}J_zu_r,\qquad r=1,2,
\]
are mutually orthonormal. Hence, with $\Phi(X)=\sum_aJ_aXJ_a$,
\[
 Q=I-P-\frac3{20}\Phi(P)
\]
is an orthogonal rank-one projector and $PQ=0$.

\paragraph{2. Recovering $P$ from the multipoles of $Q$.}
Write $X=\sum_{L=0}^8 X_L$ for the Hilbert--Schmidt orthogonal decomposition into irreducible spin-tensor ranks. The Casimir identity gives
\[
 \Phi(X_L)=\left(20-\frac{L(L+1)}2\right)X_L.
\]
Consequently, for $L>0$,
\[
 Q_L=\left(\frac{3L(L+1)}{40}-4\right)P_L,
 \qquad P_L=\beta_L Q_L,
 \qquad \beta_L=\frac{40}{3L(L+1)-160}.
\]
All denominators are nonzero for $1\le L\le8$. At rank zero,
$P_0=2I/9$ and $Q_0=I/9$, so set $\beta_0=2$.
Since $P_1=P_2=0$, also $Q_1=Q_2=0$: the complementary pure state is second-order anticoherent.
Define
\[
 a_L=\|Q_L\|_{\HS}^2\ge0,\qquad a_0=\frac19.
\]
Then
\[
 a_1=a_2=0,\qquad
 X:=\Tr(PQ)=\sum_{L=0}^8\beta_La_L=0,\qquad
 Y:=\Tr(P^2)=\sum_{L=0}^8\beta_L^2a_L=2.
\]

\paragraph{3. An exact two-copy identity.}
Let $Q=|q\rangle\langle q|$ and let $p_J$ be the total-spin-$J$ probability in $|q\rangle\otimes|q\rangle$. Only $J=0,2,4,6,8$ occur, since the two copies are identical. For Hilbert--Schmidt orthonormal spin tensors,
\[
 a_L=\sum_{J\in\{0,2,4,6,8\}}F_{LJ}p_J,
 \qquad
 F_{LJ}=(2L+1)
 \begin{Bmatrix}4&4&J\\4&4&L\end{Bmatrix},
 \qquad \sum_Jp_J=1.
\]
This can equivalently be obtained by evaluating
$\sum_M T_{LM}\otimes T_{LM}^\dagger$ on a highest-weight vector of each total-spin sector. The accompanying verifier computes all 45 coefficients by rational highest-weight sums, without invoking a Wigner or Clebsch--Gordan library.

Substitution of these rational coefficients gives the following identity for every pure $Q$ and the operator $P$ reconstructed above:
\begin{align*}
20925216a_5={}&10544114
 +\frac{5610400200}{121}a_1
 -\frac{315270952920}{5041}a_2\\
 &+25055422X-5551945Y.
\end{align*}
It is an identity in each of the five total-spin sectors, not a numerical inequality. The complete rational matrix and all five coefficient checks accompany this proof.
For the hypothetical subspace, $a_1=a_2=X=0$ and $Y=2$. Dividing the resulting equality by 96 yields
\[
 \boxed{\quad 0=5831+217971a_5.\quad}
\]
The right-hand side is strictly positive because $a_5\ge0$. This contradiction excludes every two-dimensional second-order anticoherent subspace.

\paragraph{4. Attainment and third order.}
The normalized state
\[
 |\psi\rangle=\sqrt{\frac5{24}}
 \bigl(|D_8^{(0)}\rangle+|D_8^{(8)}\rangle\bigr)
 -\sqrt{\frac7{12}}|D_8^{(4)}\rangle
\]
is third-order anticoherent. Its Dicke support gaps are four, so all off-diagonal moments through rank three vanish. Symmetry makes the odd diagonal moments vanish, and
$\langle J_z^2\rangle=20/3$. Alternatively, its complete two- and three-qubit reduced matrices are checked exactly by Dicke splitting in the verifier.
Thus dimension one is attained at both orders. A third-order subspace is also second-order, so the exclusion proves both claimed maxima.
\end{document}
